% Standalone insert for the revised gamma paper.
% Expected packages: amsmath, amssymb, amsthm, mathtools.
% Suggested bibliography keys are listed at the end of this file.

\section{Geronimus--Markov reduction and the Laguerre--Kummer tower}
\label{sec:geronimus-kummer}

This section replaces finite symbolic verification by an all-degree argument.
It also separates three different statements that should not be conflated:
the exact determinant identity, the classical Painlev\'e specialization, and
the arithmetic usefulness of the resulting linear forms.

Let \(d\nu\) be a positive measure of infinite support on
\([0,\infty)\), with finite moments
\[
  m_k=\int_0^\infty x^k\,d\nu(x) \qquad(k\geq 0).
\]
Write \(p_j\) for its monic orthogonal polynomials and
\[
  \Delta_j=\det(m_{r+c})_{0\leq r,c<j},\qquad \Delta_0=1.
\]
For \(s>0\), set
\[
  \mu_k(s)=\int_0^\infty\frac{x^k}{x+s}\,d\nu(x),
  \qquad
  \mathcal D_n(s)=\det(\mu_{r+c}(s))_{0\leq r,c<n}.
\]

\begin{theorem}[Geronimus--Markov determinant reduction]
\label{thm:GM-reduction}
For every \(n\geq1\),
\begin{equation}
 \boxed{
 \mathcal D_n(s)=(-1)^{n-1}\Delta_{n-1}
 \int_0^\infty\frac{p_{n-1}(x)}{x+s}\,d\nu(x).}
 \label{eq:GM-reduction}
\end{equation}
If \(F(s)=\mu_0(s)\), then equivalently
\begin{equation}
 \begin{aligned}
 \mathcal D_n(s)
   &=\Delta_{n-1}\bigl(Q_{n-1}(s)F(s)-P_{n-2}(s)\bigr),\\
 Q_{n-1}(s)&=(-1)^{n-1}p_{n-1}(-s)>0,
 \end{aligned}
 \label{eq:Markov-Pade}
\end{equation}
where \(P_{n-2}\) is a polynomial of degree at most \(n-2\), with the
convention \(P_{-1}=0\).
Thus the Cauchy-transform Hankel determinant is precisely the classical
Markov--Pad\'e remainder multiplied by the preceding moment determinant.
\end{theorem}

\begin{proof}
In the defining matrix for \(\mathcal D_n\), perform the column operations
\[
 C_j\longleftarrow C_j+sC_{j-1},
 \qquad j=n-1,n-2,\ldots,1,
\]
in this order.  Since
\(\mu_k+s\mu_{k-1}=m_{k-1}\) for \(k\geq1\), the resulting columns are
\[
  (\mu_i)_{i=0}^{n-1},
  (m_i)_{i=0}^{n-1},
  \ldots,
  (m_{i+n-2})_{i=0}^{n-1}.
\]
Expansion in the first column is the functional
\(f\mapsto\int f(x)(x+s)^{-1}d\nu(x)\) applied to a polynomial orthogonal,
with respect to \(d\nu\), to \(1,x,\ldots,x^{n-2}\).  Its leading coefficient
is the cofactor \((-1)^{n-1}\Delta_{n-1}\).  The polynomial is therefore
\((-1)^{n-1}\Delta_{n-1}p_{n-1}\), proving
\eqref{eq:GM-reduction}.

Now write
\[
 \frac{p_{n-1}(x)}{x+s}
 =\frac{p_{n-1}(-s)}{x+s}
  +\frac{p_{n-1}(x)-p_{n-1}(-s)}{x+s}.
\]
The second quotient is a polynomial of degree at most \(n-2\), which gives
\eqref{eq:Markov-Pade}.  Finally, all zeros of \(p_{n-1}\) lie in
\((0,\infty)\), so \((-1)^{n-1}p_{n-1}(-s)>0\).
\end{proof}

\begin{remark}[Necessary scope]
For an arbitrary signed or finitely supported weight, the rank-one argument
only proves degree at most one in \(F\).  Exact degree one follows under the
positive infinite-support assumptions above.  Moreover,
\eqref{eq:Markov-Pade} is a structural Pad\'e identity; it does not by itself
compare arithmetic heights.  Any assertion that the determinant is
arithmetically weaker than Pad\'e requires a separate estimate for
\(\Delta_{n-1}\).
\end{remark}

We now specialize to the Laguerre measure.  For \(a>-1\), put
\begin{align*}
 d\nu_a(x)&=x^ae^{-x}\,dx,\\
 D_n^{(a)}(s)&=
 \det\left(\int_0^\infty
 \frac{x^{i+j+a}e^{-x}}{x+s}\,dx\right)_{0\leq i,j<n},\\
 C_n(a)&=\prod_{k=0}^{n-1}k!\,\Gamma(k+a+1)
 =\frac{G(n+1)G(n+a+1)}{G(a+1)}.
\end{align*}
Here and below \(G\) denotes the Barnes \(G\)-function.

\begin{theorem}[Exact Laguerre--Kummer collapse]
\label{thm:Laguerre-Kummer-exact}
For every \(n\geq1\), \(a>-1\), and \(s>0\),
\begin{equation}
 \boxed{D_n^{(a)}(s)=C_n(a)U(n,1-a,s).}
 \label{eq:Laguerre-Kummer-exact}
\end{equation}
Consequently \(y=D_n^{(a)}\) satisfies
\begin{equation}
 sy''+(1-a-s)y'-ny=0.
 \label{eq:Kummer-Dn}
\end{equation}
If
\(\Sigma_n=s(\log D_n^{(a)})'\), then
\begin{equation}
 s\Sigma_n'=(s+a)\Sigma_n-\Sigma_n^2+ns.
 \label{eq:Riccati-Sigma}
\end{equation}
In addition, if \(u_n=U(n,1-a,s)\) and \(u_0=1\), then
\begin{equation}
 n(n+a)u_{n+1}=(s+2n+a-1)u_n-u_{n-1}
 \qquad(n\geq1).
 \label{eq:Kummer-discrete}
\end{equation}
\end{theorem}

\begin{proof}
The monic Laguerre polynomial is
\(p_{n-1}(x)=(-1)^{n-1}(n-1)!L_{n-1}^{(a)}(x)\), and
\(\Delta_{n-1}=C_{n-1}(a)\).  The Laplace representation of
\((x+s)^{-1}\) and the Laguerre transform give
\begin{equation}
 \int_0^\infty x^ae^{-(1+t)x}L_{n-1}^{(a)}(x)\,dx
 =\frac{\Gamma(n+a)}{(n-1)!}
   \frac{t^{n-1}}{(1+t)^{n+a}}.
 \label{eq:Laguerre-Laplace}
\end{equation}
Substitution in \eqref{eq:GM-reduction}, followed by
\begin{equation}
 U(n,1-a,s)=\frac1{\Gamma(n)}\int_0^\infty
 e^{-st}\frac{t^{n-1}}{(1+t)^{n+a}}\,dt,
 \label{eq:U-integral}
\end{equation}
yields \eqref{eq:Laguerre-Kummer-exact}, because
\[
 C_{n-1}(a)\Gamma(n)\Gamma(n+a)=C_n(a).
\]
Equations \eqref{eq:Kummer-Dn} and \eqref{eq:Riccati-Sigma} follow from
Kummer's equation.  Equation \eqref{eq:Kummer-discrete} is the monic
Laguerre three-term recurrence for its function of the second kind.  It can
also be checked directly from \eqref{eq:U-integral}: integrate the derivative
of
\[
 e^{-st}t^{n-1}(1+t)^{-n-a}.
\]
For \(n\geq2\) both endpoint terms vanish and elementary rearrangement gives
\eqref{eq:Kummer-discrete}; for \(n=1\), the endpoint at zero is \(-1\), which
is precisely the term \(-u_0\).
\end{proof}

\begin{remark}[Literature status]
The mathematical content of \eqref{eq:Laguerre-Kummer-exact} is classical:
the Stieltjes transform of a monic Laguerre polynomial is its function of the
second kind and is a Kummer \(U\)-function.  See, for example,
Dea\~no--Huertas--Rom\'an~\cite{DeanoHuertasRoman2018}, especially their
equations~(4)--(5), and DLMF~\cite{DLMF}, \S13.4.  The point here is to provide
a complete determinant proof and to fix the normalization, not to claim a new
special-function identity.
\end{remark}

\subsection{Rigorous growth and the classical Pad\'e remainder}

\begin{theorem}[Superfactorial divergence]
\label{thm:Laguerre-divergence}
For fixed \(a>-1\) and \(s>0\),
\begin{equation}
 \frac{e^{-2s}}{3^{a+1}2^{n-1}\Gamma(n)}
 \leq U(n,1-a,s)\leq\frac1{s\Gamma(n)}.
 \label{eq:U-bounds}
\end{equation}
Consequently
\begin{equation}
 \boxed{
 \log D_n^{(a)}(s)=n^2\log n-\frac32n^2+O(n\log n),}
 \qquad
 D_n^{(a)}(s)\longrightarrow+\infty.
 \label{eq:Dn-growth}
\end{equation}
\end{theorem}

\begin{proof}
For the lower bound in \eqref{eq:U-bounds}, restrict
\eqref{eq:U-integral} to \(1\leq t\leq2\) and write its integrand as
\[
 e^{-st}\left(\frac{t}{1+t}\right)^{n-1}(1+t)^{-a-1}.
\]
On this interval the three factors are bounded below by
\(e^{-2s}\), \(2^{-(n-1)}\), and \(3^{-a-1}\), respectively.  For the upper
bound, the last two factors have product at most one, and
\(\int_0^\infty e^{-st}dt=s^{-1}\).
Using
\[
 C_n(a)=\frac{G(n+1)G(n+a+1)}{G(a+1)}
\]
and the standard Barnes-\(G\) asymptotic now gives
\eqref{eq:Dn-growth}; see DLMF~\cite[\S5.17]{DLMF}.
\end{proof}

For \(a=0\), let \(B(s)=e^sE_1(s)=U(1,1,s)\), so that
\(B(1)=\delta\), and put
\[
 Q_{n-1}(s)=(n-1)!L_{n-1}(-s).
\]
Theorem~\ref{thm:GM-reduction} gives a polynomial
\(P_{n-2}(s)\in\mathbb Z[s]\) such that
\begin{equation}
 \boxed{
 Q_{n-1}(s)B(s)-P_{n-2}(s)=(n-1)!^2U(n,1,s).}
 \label{eq:classical-Pade-remainder}
\end{equation}
Indeed, \((n-1)!L_{n-1}(x)\in\mathbb Z[x]\), polynomial division by
\(x+s\) stays in \(\mathbb Z[x,s]\), and the moments of \(e^{-x}dx\) are
factorials.  At \(s=1\), \eqref{eq:classical-Pade-remainder} is the classical
Laguerre--Stieltjes Pad\'e remainder for the Euler--Gompertz constant.  In
particular,
\begin{equation}
 Q_{n-1}(1)\delta-P_{n-2}(1)
 \geq \frac{e^{-2}}{3}\frac{(n-1)!}{2^{n-1}}
 \longrightarrow+\infty.
 \label{eq:Pade-divergence}
\end{equation}
Thus both the determinant and its underlying classical Pad\'e linear form
diverge; this is a theorem, not numerical evidence.  The continued-fraction
identification is classical; see
Lagarias~\cite[\S2.5]{Lagarias2013} and
Ernvall-Hyt\"onen--Matala-aho--Sepp\"al\"a~\cite{EMHS2023}.

\subsection{Hausdorff moments, total positivity, and random matrices}

\begin{proposition}[Complete monotonicity and Tur\'an inequality]
\label{prop:U-Hausdorff}
Define
\[
 h_m(a,s)=m!U(m+1,1-a,s)\qquad(m\geq0).
\]
Then
\begin{equation}
 \boxed{
 h_m(a,s)=\int_0^1 r^m(1-r)^{a-1}
 \exp\left(-\frac{sr}{1-r}\right)dr.}
 \label{eq:U-Hausdorff}
\end{equation}
Consequently \((h_m)_{m\geq0}\) is a strict Hausdorff moment sequence.  With
\(\Delta f_m=f_{m+1}-f_m\),
\begin{equation}
 (-1)^k\Delta^kh_m>0\qquad(k,m\geq0).
 \label{eq:complete-monotone}
\end{equation}
It is strictly log-convex, and therefore
\begin{equation}
0<n\frac{U(n+1,1-a,s)}{U(n,1-a,s)}<1,
\qquad
n\frac{U(n+1,1-a,s)}{U(n,1-a,s)}\nearrow1,
\qquad(n\geq1).
\label{eq:U-ratio}
\end{equation}
and
\begin{equation}
 \frac{n}{n+1}U(n+1,1-a,s)^2
 <U(n,1-a,s)U(n+2,1-a,s)
 \qquad(n\geq1).
 \label{eq:U-Turan}
\end{equation}
\end{proposition}

\begin{proof}
In \eqref{eq:U-integral}, set \(r=t/(1+t)\) to obtain
\eqref{eq:U-Hausdorff}.  Formula \eqref{eq:complete-monotone} follows by
inserting the factor \((1-r)^k\) under the integral.  Strict log-convexity is
Cauchy--Schwarz, since the representing measure is not a point mass.  The
successive moment ratios are strictly increasing and tend to the essential
supremum of the support, which is one; this gives \eqref{eq:U-ratio}.
Rewriting strict log-convexity for \(h_n=n!U(n+1,1-a,s)\) gives
\eqref{eq:U-Turan}.
\end{proof}

\begin{proposition}[Total positivity and determinacy]
\label{prop:Laguerre-total-positive}
For \(a>-1\) and \(s>0\), the infinite Hankel matrix
\[
 \left(\int_0^\infty\frac{x^{i+j+a}e^{-x}}{x+s}\,dx\right)_{i,j\geq0}
\]
is strictly totally positive.  Its Stieltjes moment problem is determinate.
\end{proposition}

\begin{proof}
Andreief's identity expresses a minor with increasing row exponents
\(r_1<\cdots<r_k\) and column exponents \(c_1<\cdots<c_k\) as the integral of
\[
 \det(x_j^{r_i})_{i,j=1}^k\det(x_j^{c_i})_{i,j=1}^k
 \prod_{j=1}^k\frac{x_j^ae^{-x_j}}{x_j+s}\,dx_j.
\]
On \(0<x_1<\cdots<x_k\), each alternant is a Vandermonde determinant times
a Schur polynomial and is strictly positive.  Hence every minor is positive.
For determinacy, if \(\mu_k\) denotes the moment sequence, then
\[
 \mu_{2k}\leq s^{-1}\Gamma(2k+a+1),
\]
so \(\sum_k\mu_{2k}^{-1/(2k)}=\infty\); Carleman's criterion applies.
\end{proof}

\begin{corollary}[Laguerre-unitary-ensemble interpretation]
\label{cor:LUE-inverse-characteristic}
Let \(X\) have the \(n\)-dimensional Laguerre unitary ensemble with parameter
\(a>-1\), normalized so that its eigenvalue density is proportional to
\[
 \prod_{i<j}(x_i-x_j)^2\prod_{i=1}^n x_i^ae^{-x_i}.
\]
Then
\begin{equation}
 \boxed{\mathbb E\det(X+sI)^{-1}=U(n,1-a,s).}
 \label{eq:LUE-inverse-characteristic}
\end{equation}
\end{corollary}

\begin{proof}
Andreief's identity identifies the ratio of the pole-deformed and undeformed
partition functions with \(D_n^{(a)}(s)/C_n(a)\).  Apply
\eqref{eq:Laguerre-Kummer-exact}.
\end{proof}

Equation \eqref{eq:LUE-inverse-characteristic} is the standard average inverse
characteristic polynomial, or Laguerre function of the second kind.  It is a
useful interpretation, not a new random-matrix problem solved here.

\subsection{The exact Painlev\'e V specialization and its boundary condition}

The Riccati equation \eqref{eq:Riccati-Sigma} is not itself the standard
Jimbo--Miwa--Okamoto sigma form.  Put
\[
 H_n(s)=\Sigma_n(s)+n.
\]
We use the JMO convention~\cite[Appendix~C, Eq.~(C.45)]{JimboMiwa1981}
\[
 (sH'')^2=
 \bigl[H-sH'+2(H')^2+(\nu_0+\nu_1+\nu_2+\nu_3)H'\bigr]^2
 -4\prod_{j=0}^3(H'+\nu_j).
\]
Direct substitution gives
\begin{equation}
\begin{split}
 (sH_n'')^2={}&
 \bigl[H_n-sH_n'+2(H_n')^2-(2n+a+1)H_n'\bigr]^2\\
 &-4H_n'(H_n'-1)(H_n'-n)(H_n'-n-a).
\end{split}
\label{eq:JMO-PV}
\end{equation}
This is the JMO sigma form of \(P_V\) with
\[
 (\nu_0,\nu_1,\nu_2,\nu_3)=(0,-1,-n,-n-a).
\]
For \(a>0\), it is the \(\lambda=-1\) specialization of the deformed Laguerre
theorem of Mu--Lyu~\cite{MuLyu2025}; for \(-1<a\leq0\), the identification
follows directly from \eqref{eq:Riccati-Sigma}.  Equation
\eqref{eq:Riccati-Sigma} also shows that this specialization lies on a
classical Riccati locus.

For \(a=0\), the identity \(B'=B-s^{-1}\) is invariant under
\(B\mapsto B+ce^s\).  At the level of moments this is
\begin{equation}
 \mu_k(s)\longmapsto\mu_k(s)+ce^s(-s)^k,
 \label{eq:Uvarov-shift}
\end{equation}
namely the Uvarov mass modification
\[
 \frac{e^{-x}}{x+s}\,dx
 \longmapsto
 \frac{e^{-x}}{x+s}\,dx+ce^s\delta_{-s}.
\]
The transformed determinant is another linear combination of the two Kummer
solutions and remains on the same Riccati/Painlev\'e locus.  However, it does
not preserve the distinguished Stieltjes boundary when \(c\neq0\).  For every
\(a>-1\), that boundary is
\begin{equation}
 D_n^{(a)}(s)\sim C_n(a)s^{-n}\qquad(s\to+\infty),
 \label{eq:recessive-boundary}
\end{equation}
or equivalently
\[
 H_n(s)=\frac{n(n+a)}s+O(s^{-2}).
\]
If \(c<0\), the added atom has negative mass; if \(c\notin\mathbb R\), the
functional is not a positive real measure; and if \(c>0\), positive mass is
placed at \(-s\notin[0,\infty)\).  Thus no \(c\neq0\) preserves the Stieltjes
class under discussion.
Thus the differential equation without boundary data has the translation
symmetry, whereas the actual Stieltjes and random-matrix solution is
canonically selected.  No arithmetic ``no-go'' statement follows from the
formal shift alone.  The direct \(B\)-shift applies to \(a=0\), and after
reduction to \(B\) to nonnegative integral \(a\); it is not an action on the
arbitrary-real-\(a\) family.

\section{A confluent reciprocal-Gamma matching determinant}
\label{sec:reciprocal-gamma-matching}

Put \(H_0=0\), \(H_m=\sum_{r=1}^m r^{-1}\), and
\[
 K_N(x)=\det\left(\frac{H_{i+j}-x}{(i+j)!}\right)_{0\leq i,j\leq N},
 \qquad
 M_N=\binom{N+1}{2}.
\]
Let \(\mathfrak M_N\) be the set of partial matchings of
\(\{0,1,\ldots,N\}\), and write \(V(\mathfrak m)\) for the matched vertices.
Define
\begin{align*}
 h_{N,j}&=H_{N+j}-H_j+H_{N-j},\\
 u_{N,j}(x)&=x-h_{N,j},\\
 B_N&=\frac{\prod_{k=1}^Nk!}{\prod_{m=N}^{2N}m!}.
\end{align*}

\begin{theorem}[Reciprocal-Gamma matching formula]
\label{thm:reciprocal-gamma-matching}
For every \(N\geq1\),
\begin{equation}
 \boxed{
 K_N(x)=(-1)^{N+1+M_N}B_N Z_N(x),}
 \label{eq:KN-matching-prefactor}
\end{equation}
where
\begin{equation}
 Z_N(x)=
 \sum_{\mathfrak m\in\mathfrak M_N}
 \left(\prod_{\{j,k\}\in\mathfrak m}\frac1{(j-k)^2}\right)
 \left(\prod_{j\notin V(\mathfrak m)}u_{N,j}(x)\right).
 \label{eq:ZN-matching}
\end{equation}
In particular \(\deg K_N=N+1\), with leading coefficient
\((-1)^{N+1+M_N}B_N\).
\end{theorem}

\begin{proof}
Introduce one auxiliary variable per column and set
\[
 F_N(\boldsymbol\alpha)=
 \det\left(\frac1{\Gamma(i+j+1+\alpha_j)}\right)_{0\leq i,j\leq N}.
\]
A reciprocal-Gamma Vandermonde evaluation gives
\begin{equation}
 F_N(\boldsymbol\alpha)=(-1)^{M_N}
 \frac{\prod_{0\leq j<k\leq N}
 (k-j+\alpha_k-\alpha_j)}
 {\prod_{j=0}^N\Gamma(N+j+1+\alpha_j)}.
 \label{eq:reciprocal-gamma-Vandermonde}
\end{equation}
Since
\[
 -\left.(\partial_{\alpha_j}+x-\gamma)
 \frac1{\Gamma(i+j+1+\alpha_j)}\right|_{\alpha_j=0}
 =\frac{H_{i+j}-x}{(i+j)!},
\]
we apply one such operator in every column.  At
\(\boldsymbol\alpha=0\), one has
\(F_N(\boldsymbol0)=(-1)^{M_N}B_N\neq0\), so a local logarithm can be chosen.
Equivalently, put
\[
 \widetilde F_N(\boldsymbol\alpha)
 =e^{(x-\gamma)\sum_j\alpha_j}F_N(\boldsymbol\alpha).
\]
Then
\[
 \left.\prod_{j=0}^N(\partial_{\alpha_j}+x-\gamma)F_N
 \right|_{\boldsymbol0}
 =
 \left.\partial_{\alpha_0}\cdots\partial_{\alpha_N}\widetilde F_N
 \right|_{\boldsymbol0}.
\]
The distinct-variable logarithmic cumulants of \(\widetilde F_N\) are
\begin{align*}
 \kappa_{\{j\}}
 &=\left.\partial_{\alpha_j}\log F_N\right|_{\boldsymbol0}+x-\gamma
   =u_{N,j}(x),\\
 \kappa_{\{j,k\}}
 &=\left.\partial_{\alpha_j}\partial_{\alpha_k}\log F_N
   \right|_{\boldsymbol0}=(j-k)^{-2},\\
 \kappa_S
 &=\left.\prod_{j\in S}\partial_{\alpha_j}\log F_N
   \right|_{\boldsymbol0}=0\qquad(|S|\geq3).
\end{align*}
The multivariate moment--cumulant formula therefore sums precisely over set
partitions into singletons and pairs, which are the partial matchings in
\eqref{eq:ZN-matching}.  Evaluation of \(F_N(\boldsymbol0)\), together with
the \(N+1\) minus signs, gives
\eqref{eq:KN-matching-prefactor}.
\end{proof}

The matching sum has two exact cross-field interpretations.  First, it is the
monomer--dimer partition function of the complete graph on
\(\{0,\ldots,N\}\), with monomer activities \(u_{N,j}(x)\) and positive dimer
activities \((j-k)^{-2}\).  Second, let \(A_N\) be a random real
skew-symmetric matrix whose upper-triangular entries are independent and
centered, with
\[
 \mathbb E(A_N)_{jk}^2=(j-k)^{-2}\qquad(j<k),
\]
and let
\(\mathsf H_N=\operatorname{diag}(h_{N,0},\ldots,h_{N,N})\).  Expansion by
permutations gives
\begin{equation}
 \boxed{Z_N(x)=\mathbb E\det(xI-\mathsf H_N+A_N).}
 \label{eq:skew-RMT-matching}
\end{equation}
Indeed, every cycle of length at least three contains independent centered
entries to the first power and has zero expectation, while the transpositions
give the dimer factors.  Formula \eqref{eq:skew-RMT-matching} is an expected
characteristic-polynomial representation, not a random-matrix limit theorem.

\begin{theorem}[Heilmann--Lieb complex-zero strip]
\label{thm:KN-zero-strip}
Every complex zero \(\zeta\) of \(K_N\) satisfies
\begin{equation}
 \boxed{H_{2N}-H_N\leq\Re\zeta\leq2H_N.}
 \label{eq:KN-zero-strip}
\end{equation}
\end{theorem}

\begin{proof}
A direct difference calculation shows that \(j\mapsto h_{N,j}\) is strictly
decreasing, with
\[
 h_{N,0}=2H_N,
 \qquad
 h_{N,N}=H_{2N}-H_N.
\]
If \(\Re x>2H_N\), every monomer activity
\(u_{N,j}(x)=x-h_{N,j}\) lies in the open right half-plane.  If
\(\Re x<H_{2N}-H_N\), they all lie in the open left half-plane.  The
Heilmann--Lieb left--right half-plane theorem for a monomer--dimer partition
function with nonnegative dimer activities~\cite[Theorem~4.6 and
Lemma~4.7]{HeilmannLieb1972} says that \(Z_N(x)\neq0\) in either case; see also
the precise multivariate restatement in
\cite[Theorem~10.1]{ChoeOxleySokalWagner2004}.  Equation
\eqref{eq:KN-matching-prefactor} proves the claim.
\end{proof}

\begin{remark}
Theorem~\ref{thm:KN-zero-strip} does not assert real-rootedness.  In fact
\[
 K_1(x)=-\frac{2x^2-5x+4}{4}
\]
already has nonreal roots.  Heilmann--Lieb supplies a zero-free complement of
a vertical strip, not a real-rooted matching polynomial.
\end{remark}

\begin{proposition}[Cleared determinants below the strip]
\label{prop:KN-divergence}
Fix a real \(x<\log2\).  If \(q_N\) is any positive integer that clears all
coefficient denominators of \(K_N\), then, for all sufficiently large \(N\),
\begin{equation}
 q_N|K_N(x)|
 \geq\prod_{j=0}^N(h_{N,j}-x)
 \longrightarrow+\infty.
 \label{eq:KN-cleared-divergence}
\end{equation}
\end{proposition}

\begin{proof}
For large \(N\), \(x<h_{N,N}=H_{2N}-H_N\).  Hence all singleton factors
\(u_{N,j}(x)\) are negative.  A matching with \(r\) pairs leaves
\(N+1-2r\) singletons, so every term in \eqref{eq:ZN-matching} has the same
sign \((-1)^{N+1}\).  There is no cancellation, and the empty matching gives
\[
 |Z_N(x)|\geq\prod_{j=0}^N(h_{N,j}-x).
\]
The leading coefficient of \(K_N\) is \(\pm B_N\).  Since \(q_NK_N\) has
integer coefficients, \(q_NB_N\) is a positive integer and is at least one.
This proves the inequality.  Finally, for \(j\leq N-1\),
\[
 h_{N,j}\geq h_{N,N-1}
 =H_{2N-1}-H_{N-1}+1\longrightarrow1+\log2,
\]
whereas \(h_{N,N}-x\to\log2-x>0\).  At least \(N\) factors are therefore
eventually bounded below by a fixed number greater than one, proving
divergence.
\end{proof}

In particular, \eqref{eq:KN-cleared-divergence} applies at \(x=\gamma\): an
elementary estimate gives \(\gamma<\log2\).  For completeness, using
\[
 \gamma=\sum_{k\geq1}\left(\frac1k-\log\left(1+\frac1k\right)\right),
 \qquad
 u-\log(1+u)<\frac{u^2}{2},
\]
one obtains \(\gamma<11/8-\log2<\log2\); the last inequality follows from
\(\log2>2(1/3+1/81)=56/81>11/16\).  Hence
\[
 q_N|K_N(\gamma)|\longrightarrow+\infty.
\]
This is an obstruction for this cleared-determinant normalization, not an
irrationality or transcendence statement about \(\gamma\).

\begin{remark}[Priority and scope]
The Geronimus reduction and the Laguerre function of the second kind are
classical.  So are the continued fraction, the Uvarov modification, and the
Painlev\'e V embedding, or they are direct specializations of cited results.
The zero strip is a direct application of the classical Heilmann--Lieb theorem
to the exact matching formula.  These statements give rigorous structural
closures.  They also rule out the specified determinant normalizations as
sources of small linear forms.
They do not solve a named open problem in orthogonal polynomials, moment
theory, continued fractions, isomonodromy, or random matrices.  The exact
priority of the differentiated reciprocal-Gamma matching formula itself
requires a specialist literature search before any novelty claim.
\end{remark}

% Suggested bibliography entries/keys for the main manuscript:
% DeanoHuertasRoman2018:
% A. Deano, E. J. Huertas, P. Roman, Asymptotics of orthogonal polynomials
% generated by a Geronimus perturbation of the Laguerre measure,
% J. Math. Anal. Appl. 433 (2016), 732--746; arXiv:1504.05976.
%
% MuLyu2025:
% X. Mu, S. Lyu, Hankel determinants for a deformed Laguerre weight with
% multiple variables and generalized Painleve V equation,
% arXiv:2410.22780; J. Math. Phys. 66 (2025), 073506.
%
% JimboMiwa1981:
% M. Jimbo, T. Miwa, Monodromy preserving deformation of linear ordinary
% differential equations with rational coefficients. II,
% Physica D 2 (1981), 407--448.
%
% HeilmannLieb1972:
% O. J. Heilmann, E. H. Lieb, Theory of monomer-dimer systems,
% Comm. Math. Phys. 25 (1972), 190--232.
%
% ChoeOxleySokalWagner2004:
% Y.-B. Choe, J. G. Oxley, A. D. Sokal, D. G. Wagner,
% Homogeneous multivariate polynomials with the half-plane property,
% Adv. in Appl. Math. 32 (2004), 88--187.
%
% Lagarias2013:
% J. C. Lagarias, Euler's constant: Euler's work and modern developments,
% Bull. Amer. Math. Soc. 50 (2013), 527--628; arXiv:1303.1856.
%
% EMHS2023:
% A.-M. Ernvall-Hytonen, T. Matala-aho, L. Seppala,
% Euler's factorial series, Hardy integral, and continued fractions,
% J. Number Theory 244 (2023), 224--250; arXiv:2111.13649.
%
% DLMF:
% NIST Digital Library of Mathematical Functions, Chapters 5, 13, and 18.
